Conditional prediction can score individual cells of a dirty table, but the
reference examples used by a tabular foundation model can themselves be corrupted.
We evaluate a fixed purification rule that removes context rows with the largest
nested out-of-fold target-cell scores, with all first-pass work isolated from the
outer query fold. A paired random-removal control holds base contexts and removal
counts fixed, and nested preparation costs are reported explicitly. The study covers
four capped benchmark tables, real and two injected-error conditions, five seeds,
and both TabPFN-2 and LightGBM. On real errors, purification was worse than matched
random removal on three tables and better on one (TabPFN AP difference $-0.161$
on hospital, $-0.049$ on flights, $+0.026$ on beers and $-0.030$ on rayyan), and no detector,
including the simple baselines, was best on every table. These are per-table
findings with conditional row-bootstrap uncertainty, not a population-wide verdict.
We also report support, failures, rare-valid alarms, boundary sensitivity and
compute on the complete eligible-cell universe. The results argue for testing
context selection against matched random removal rather than assuming that a
cleaner-looking context helps.
